% =====================================================================
%  Chapter 14 — Profile rigidity (S39)
% =====================================================================
\chapter{Profile Rigidity, and the Ceiling}
\label{ch:profile-rigidity}

\begin{record}
\textbf{What this chapter covers.} Session \Sn{39}, and the section
\texttt{sec:rigidity-s39} it produced --- the largest single addition the proof
document ever received (\(6667\to7737\) lines). Its idea is a change of
object: instead of trying to force \(\sstar(t)\to0\) as a statement about a
solution's behaviour in time, pass to a Type-I blow-up profile and ask whether
the \emph{limit object} can have a given value of \(\sstar\). Decay becomes
rigidity. \(\sstar\) stops being a quantity that must be shown to fall and
becomes an invariant of an ancient solution.

\smallskip
This chapter also carries the sentence that bounds the whole later programme.
\texttt{rem:s39-distance} (line~7703) states that even with every remaining
hypothesis discharged and both horns closed, the conclusion is \emph{no
Type-I singularity} --- ``strictly less than the Clay problem''. That is the
ceiling. It has not moved since, including at \Sn{48}, and this chapter states
it plainly rather than at the end.
\end{record}

\section{Why the object had to change}
\label{sec:why-reframe}

By the close of \Sn{38} the position was exact and unpromising. One deep open
problem remained --- force \(\sstar(t)\to0\) along a putative Type-I blow-up
sequence --- and \Sn{38} had proved it \emph{necessary}, closing the only
alternative. Every direct attempt on it had failed in the same way, three
times, over the programme's history:

\begin{itemize}[leftmargin=1.6em]
\item \S19's gradient recursion produced a sublinear exponent
  (\texttt{rem:prop19-gap}) and reduced to CKN;
\item \S20's vorticity recursion produced the \Sn{31} audit;
\item \Sn{38}'s own closure target --- assembling the proved first rung,
  second rung and amended \(K\)-absorption into a scale recursion --- could
  not be shown subcritical, and was deliberately left \OPEN{} rather than
  forced.
\end{itemize}

All three are the same manoeuvre: iterate an inequality down scales and hope
the effective coefficient falls below one. \Sn{39} stopped attempting it.

\section{The construction}
\label{sec:rigidity-construction}

\subsection{Rescaling, and the ancient class}

Fix a putative Type-I singularity at \((x^*,T)\) and set
\begin{equation}\label{eq:rescale}
  u_\lambda(y,s) \;=\; \lambda\,u\!\left(x^*+\lambda y,\;T+\lambda^2 s\right),
  \qquad \lambda\in(0,1],\; s<0 .
\end{equation}
Under the Type-I bounds --- both of them, the vorticity bound
\(M(t)\le C_I/(T-t)\) \emph{and} the velocity bound
\(\sup\abs{u}\le C_I(T-t)^{-1/2}\), which line~5565 of the proof document
assumes outright --- the rescaled family is uniformly bounded, and a
subsequence converges to an \emph{ancient} solution \(U\) on
\(\R^3\times(-\infty,0)\) in the class \(\mathcal{A}(c_I)\).

\texttt{lem:typeI-scaleinv-bounds} (line~6729) proves that this class is
closed under the Navier--Stokes scaling group \emph{with the same constant},
and that \(\sstar(U_\mu)(s)=\sstar(U)(\mu^2 s)\). That invariance is not
decoration: it is what makes a \emph{second} compactness extraction legal
later, in \texttt{prop:two-step-reduction}.

\begin{figure}[ht]
\centering
\begin{tikzpicture}[
  >={Latex[length=2mm]},
  bx/.style={draw, rounded corners=2pt, font=\scriptsize, align=center,
             inner sep=3.5pt, text width=34mm, minimum height=10mm},
  hy/.style={bx, draw=openblue, thick, fill=panelgrey},
  gd/.style={bx, draw=provedgreen, very thick},
  op/.style={bx, draw=openblue, thick},
  ar/.style={->, thick}
]
\node[bx] (u) at (0,0)
  {\(u\) on \(\T^3\times[0,T)\)\\ Type-I: \emph{both} bounds\\ (line~5565)};
\node[hy] (fam) at (4.8,0)
  {rescaled family \eqref{eq:rescale}\\
   uniform in \(\lambda\) by \textbf{(R)}};
\node[gd] (U) at (9.6,0)
  {ancient profile \(U\)\\ on \(\R^3\times(-\infty,0)\)\\
   \(\sstar(U)\) an \emph{invariant}};

\draw[ar] (u) -- node[above, font=\scriptsize] {\(\lambda\to0\)} (fam);
\draw[ar] (fam) -- node[above, font=\scriptsize] {\(C^\infty_{\mathrm{loc}}\)} (U);
\node[font=\scriptsize, below=1mm of fam, text=rulegrey, text width=34mm,
      align=center] {needs \textbf{(N)}: the strength must not escape the
      window};

\node[gd] (horn1) at (6.6,-3.0)
  {\textbf{Aligned horn}\\ \(\sstar(U)\equiv0\)\\
   \(\Rightarrow\) \(\partial_3U=0\), 2D\\
   \(\Rightarrow\) Liouville \(\Rightarrow\) trivial};
\node[op] (horn2) at (12.0,-3.0)
  {\textbf{Disordered horn}\\ \(\sstar(U)\ge\delta>0\)\\
   over an infinite past\\
   \OPEN{}: \texttt{target:disorder-depletion}};

\draw[ar] (U.south) to[out=-120,in=90] (horn1.north);
\draw[ar] (U.south) to[out=-60,in=90] (horn2.north);

\node[bx, draw=impossiblered, thick, text width=52mm] (ceiling) at (5.0,-5.6)
  {\textbf{Ceiling} (\texttt{rem:s39-distance}, line~7703):
   both horns closed \(\Rightarrow\) \emph{no Type-I singularity} ---
   strictly less than the Clay problem};
\draw[->, thick, impossiblered, dashed] (horn1.south) -- (ceiling.east);
\end{tikzpicture}
\caption[The profile-rigidity reframe]{The \Sn{39} construction. A putative
Type-I singularity is rescaled to an ancient limit profile on which
\(\sstar\) is a scale invariant rather than a decaying quantity; the argument
then splits into two horns. The dashed arrow is the part that must be kept in
view: closing both horns gives a Type-I exclusion, not the Prize.}
\label{fig:rigidity}
\end{figure}

\subsection{Step 1: extraction, and the two hypotheses it costs}

\texttt{prop:profile-extraction} (line~6870) performs the extraction, and it
carries two hypotheses that the rest of the chapter has to keep track of.

\begin{statement}[\textbf{(R)}, scale-uniform interior regularity]
\label{st:R}
\(\norm{\nabla^k u_\lambda}_{L^\infty}\le C_k(c_I)(1+R)S_1^{-(k+1)/2}\),
uniformly in \(\lambda\), which upgrades the convergence to
\(C^\infty_{\mathrm{loc}}\).
Status at \Sn{39}: \CONDITIONAL{Serrin-type}. \textbf{Status now:} \PROVED{},
in a stronger form, by \texttt{thm:s48-R} --- see \S\ref{sec:s48-effect}.
\end{statement}

\begin{statement}[\textbf{(N)}, non-degeneracy]
\label{st:N}
The singularity's strength does not escape the rescaling window, so the limit
\(U\) is nontrivial. Status: \OPEN{}, assumed. CKN
\(\varepsilon\)-regularity supplies the right scale-invariant lower bound only
on \(Q_1(0)\), a cylinder that touches the terminal slice \(s=0\) where
compactness fails; pushing mass off it needs a Type-I-improved
\(\varepsilon\)-regularity theorem which is not proved here
(\texttt{rem:N-discussion}, line~6842).
\end{statement}

\subsection{Step 2: the transfer audit}

\texttt{prop:transfer-audit} (line~6955) asks, part by part, which of the
programme's machinery survives the limit. Its framing matters and is easy to
misread: parts (a)--(c) are \emph{re-derived from scratch on \(U\)}, not
transferred as limits of inequalities. \(U\) is a smooth solution with
scale-invariant Type-I bounds, and every one of those proofs is local and
pointwise, so they simply hold.

\begin{enumerate}[label=(\alph*),leftmargin=2.2em]
\item[(a)--(c)] The magnitude equation, the direction equation, the orthogonal
  split and Gram negativity all hold on any profile. \PROVED{}.
\item[(d)] \emph{Profiles are \(\Lfac\)-free}: on the profile ``the
  logarithmic factor \(\Lfac\) may be replaced by an absolute constant''. At
  \Sn{39} this was \PROVEDMOD{(R)}; it is now \PROVED{}
  (\texttt{cor:s48-transfer-d}, line~10257).
\item[(e)] \(\nu\langle\sstar\rangle\le C(c_I)\) on \(U\). At \Sn{39},
  \PROVEDMOD{(R)}; now \PROVEDMOD{saturation} --- see
  \S\ref{sec:s48-effect}.
\end{enumerate}

Part~(d) is the reason the S39 route looked like an escape from the
Calder\'on--Zygmund logarithm. Chapter~\ref{ch:cz-log} follows what happened
to that reading; the short version is that it was recorded at \Sn{47} as a
relocation rather than a removal (Wall~2), and \Sn{48} withdrew that reading
--- the profile route never incurs the logarithm at all, because it never
passes through the Calder\'on--Zygmund operator.

\subsection{What does not transfer}

Two prohibitions were found in the same step, and both constrain earlier work.

\begin{statement}[\texttt{rem:s39-nontransfer}(i), line~7045]
\label{st:nontransfer-i}
\IMPOSSIBLE{} as written. \texttt{prop:localized-cf} --- the entire H2 /
Regime-B engine of \Sn{37} --- does \emph{not} survive passage to the profile.
Its far-field step bounds
\(\int_{\abs{y}\ge\rho'}\abs{\omega(x^*-y)}\abs{y}^{-3}\dee y\) by
\(C(\rho')^{-3/2}\Omega^{1/2}\), and for \(U\in\mathcal{A}(c_I)\) on \(\R^3\)
there is no reason for \(\norm{\omega_U}_{L^2(\R^3)}\) to be finite; the naive
substitute \(\abs{\omega_U}\le M_U\) makes the integral logarithmically
divergent at large \(\abs{y}\). The document's phrasing: ``the entire
H2/Regime-B analysis of \S\,\texttt{sec:cf-s37} is unavailable on the profile as
written''.
\end{statement}

\begin{statement}[\texttt{rem:s39-nontransfer}(ii), line~7045]
\label{st:nontransfer-ii}
Forbidden. The Bridge Lemma may not be pulled back to the approximating
sequence. \texttt{lem:sigma-semicontinuity} (line~7075) gives
\(\sstar(U)\le\liminf\sstar(u_{\lambda_j})\) and \emph{no reverse inequality};
so \(\sstar(U)=0\) does not give \(\sstar(u_{\lambda_j})\to0\), and the hybrid
strategy ``the profile is coherent, therefore apply H1 and H2 to the original
sequence'' is invalid. The \Sn{39} dichotomy is built to obtain its
contradiction entirely on \(U\) for exactly this reason.
\end{statement}

The semicontinuity direction that \emph{is} available is the useful one:
smallness descends to the limit, which is what
\texttt{prop:two-step-reduction} (line~7557) uses to absorb the oscillatory
middle case --- \(\liminf\sstar=0\) but \(\sstar\) never zero --- into the
aligned horn by a second compactness extraction.

\section{The aligned horn}
\label{sec:aligned-horn}

Suppose \(\sstar(U)=0\). \texttt{lem:alignment-local} (line~7164) gives local
alignment: \(\ehat\) is constant on the relevant component of the
high-vorticity set. Then comes the session's genuinely new mathematics.

\begin{statement}[\texttt{lem:harmonic-reduction}, line~7245]
\label{st:harmonic}
\PROVED{} (at \Sn{39}, modulo \textbf{(R)}; unconditionally since \Sn{48}).
If \(\omega_U\parallel e_3\) globally, then \(\partial_3U\equiv0\),
\(U_3=U_3(s)\), and the horizontal part \(U_h\) solves the two-dimensional
Navier--Stokes equations exactly.
\end{statement}

The proof is five steps and was re-verified independently before the section
was merged. \(\operatorname{div}\omega=0\) forces \(\partial_3\zeta=0\), so
\(f:=\partial_3U\) is both curl-free and divergence-free, hence componentwise
harmonic; the Type-I bounds make it bounded, so Liouville forces \(f=f(s)\);
writing \(U=f(s)y_3+\tilde U\) and using boundedness of \(U\) in \(y_3\)
forces \(f\equiv0\).

This step is what makes any Liouville theorem applicable at all. General
three-dimensional bounded-ancient Liouville is open; no version of the aligned
horn can work without a genuine dimensional reduction, and this is it. With it,
\texttt{thm:aligned-horn} (line~7310) invokes the two-dimensional
bounded-ancient Liouville theorem of \cite{KNSS2009} --- with a flagged
mild-solution class caveat --- and concludes triviality.

The horn is complete except for one hypothesis.

\begin{statement}[\textbf{(A-up)}, the global alignment upgrade]
\label{st:aup}
\OPEN{} --- the largest new gap \Sn{39} created. Local, one-time,
one-component alignment must be upgraded to global \(\omega\parallel e_0\) on
\(\R^3\times(-\infty,0)\), because the harmonic Liouville step has no ball
version. \texttt{rem:A-up-discussion} (line~7187) canvasses three mechanisms:
unique continuation for the \(w\)-inequality, which reduces to the programme's
own Kato-class problem and is judged mildly positive; backward uniqueness in
the Escauriaza--Seregin--\v{S}ver\'ak sense, which lacks the spatial decay;
and scale exhaustion, which would need \(\sstar\equiv0\) for \emph{all} \(s\),
which the dichotomy does not deliver.
\end{statement}

\section{The disordered horn}
\label{sec:disordered-horn}

Suppose instead \(\sstar(U)\ge\delta>0\) persistently. The natural first
attempt --- accumulate a budget over the infinite past until it diverges ---
was named in advance as a likely trap, and \Sn{39} converted the trap into a
theorem.

\begin{statement}[\texttt{prop:no-summable-budget}, line~7356]
\label{st:nobudget}
\IMPOSSIBLE{} The naive infinite-past budget argument cannot work. Dyadic past
shells and self-similar cylinders are commensurate, and every cylinder carries
the same dimensionless bound with no decay in \(k\) --- identically so, for an
exactly self-similar profile, by scale invariance. Hence
\(\sum_k\langle\sstar\rangle=\infty\) is \emph{forced} by the class, and
divergence of the sum carries no information.
\end{statement}

That negative result dictates the only viable framework: a bounded quantity
that is monotone in \emph{logarithmic} time. \Sn{39} found one.

\begin{statement}[\texttt{prop:logtime-necessary}, line~7459]
\label{st:logtime}
\PROVED{} (modulo attainment of the spatial supremum). Put
\(\tau=\log\abs{s}\) and \(\Phi(\tau)=\abs{s}M_U(s)\in(0,c_I]\). Evaluating
the magnitude equation at the vorticity maximum --- where \(\Delta\abs{\omega}\le0\)
by Hamilton's trick --- gives
\[
  \frac{\dee\log\Phi}{\dee\tau}
  \;\ge\; 1 + \abs{s}w_{\max} - \abs{s}\alpha_{\max}.
\]
Since \(\Phi\) is bounded over infinite \(\tau\)-length,
\[
  \limsup_{\text{log-time average}}
    \bigl[\abs{s}\alpha_{\max} - \abs{s}w_{\max}\bigr] \;\ge\; 1 .
\]
\end{statement}

The sign structure is favourable: disorder drives \(\Phi\) up, and \(\Phi\) is
bounded. What remains is to make that pay.

\begin{statement}[\texttt{target:disorder-depletion}, line~7504]
\label{st:target-dd}
\OPEN{}, and \emph{marginal}. Persistent \(\sstar\ge\delta\) must force the
log-time average to be \(\le1-c(\delta)\).
\end{statement}

\texttt{rem:target-decomposition} (line~7520) splits it into (T4a), an
average-to-point statement for \(w\) at the vorticity maximum --- natural
mechanism: weak Harnack, which plugs into the existing \(K\)-family machinery
--- and (T4b), profile depletion of \(\alpha\), which is the harder half and
which, by Statement~\ref{st:nontransfer-i}, is now known to need a genuinely
different far-field control from the one \Sn{37} built.

The word \emph{marginal} is load-bearing and the document flags it as such:
the target must beat the constant \(1\) \emph{exactly}. The programme's own
history is that marginal targets are where its errors hid. The mitigating
difference, noted but not oversold, is that this target compares two averages
by a fixed constant, rather than being a recursion needing bootstrapped
smallness.

\section{\texorpdfstring{What \Sn{48} changed, and what it did not}{What S48 changed, and what it did not}}
\label{sec:s48-effect}

\Sn{48} settled hypothesis \textbf{(R)} against the literature, in a stronger
form than \Sn{39} had stated it (\texttt{thm:s48-R},
\S\,\texttt{sec:hypR-s48}). The full statement and its route are the subject of
Chapter~\ref{ch:cz-log}. Its effect on this chapter's
ledger is mechanical and is recorded in the document as an amendment: every
row reading ``mod (R)'' should be read with the (R) deleted.

One exception is recorded, and it is not small.

\begin{obsv}[Saturation: what (e) still needs]
\label{obs:saturation}
\texttt{prop:transfer-audit}(d) states \emph{two} things. The first,
\(\norm{\nabla U}_{L^\infty}\le C_1(c_I)\abs{s}^{-1}\), is now proved. The
second, \(\norm{\nabla U(\cdot,s)}_{L^\infty}\le C(c_I)M_U(s)\), carries a
\emph{pre-existing} proviso --- ``whenever the Type-I upper bound is
saturated'', i.e.\ at times with \(M_U(s)\ge c_I^{-1}\abs{s}^{-1}\). Part~(e)
uses the second form, because it measures the gradient against \(M_U\) and
takes \(\rstar=M_U(s_0)^{-1/2}\). The two forms differ by
\(\left(\abs{s}M_U(s)\right)^{-1}\), which is \(O(1)\) only on the saturated
set.

\smallskip
So (e) moves from \PROVEDMOD{(R)} to \PROVEDMOD{saturation}, not to \PROVED{}.
Whether \(M_U(s)\gtrsim\abs{s}^{-1}\) holds for all \(s<0\) or only on a set is
not settled by \textbf{(R)} and is not settled anywhere. Under \textbf{(N)},
\(M_U(s)>0\) for every \(s<0\), but positivity is weaker than saturation. The
proviso was verified as genuinely pre-existing, not introduced by the session
that found it.
\end{obsv}

Nothing else in the ledger moves. \textbf{(N)}, \textbf{(N\('\))},
\textbf{(A-up)} and \texttt{target:disorder-depletion} are untouched, and the
ceiling stands verbatim.

\section{The ledger}
\label{sec:s39-ledger}

{\footnotesize
\begin{longtable}{@{}p{0.36\textwidth} p{0.26\textwidth} p{0.30\textwidth}@{}}
\caption[The \Sn{39} ledger, with the \Sn{48} amendment]{The
\texttt{sec:rigidity-s39} ledger as the document carries it, with the
\Sn{48} amendment applied in the third column. Twenty-six rows; four of them
are proofs that something cannot be done. The last row is the one that bounds
everything above it.}
\label{tab:s39-ledger}\\
\toprule
Item & Status at \Sn{39} & Status after \Sn{48} \\
\midrule
\endfirsthead
\toprule
Item & Status at \Sn{39} & Status after \Sn{48} \\
\midrule
\endhead
\bottomrule
\endfoot
Type-I ancient class scaling-closed
  & Proved (\texttt{lem:typeI-scaleinv-bounds}) & \PROVED \\
Rescaled bounds \(\lambda\)-uniform
  & Proved & \PROVED \\
Scale-uniform interior regularity \textbf{(R)}
  & \textbf{Conditional} (Serrin-type) & \PROVED{} (\texttt{thm:s48-R}) \\
Non-degeneracy / no escape \textbf{(N)}
  & \textbf{Assumed}; CKN gives it only on \(Q_1(0)\) & \OPEN{} (unchanged) \\
Profile extraction
  & Proved mod (R),(N) & \PROVEDMOD{N} \\
Magnitude, direction, split, Gram on \(U\)
  & \textbf{Proved} (re-derived, not transferred) & \PROVED \\
Profiles are \(\Lfac\)-free
  & Proved mod (R) & \PROVED{} (\texttt{cor:s48-transfer-d}) \\
\(\nu\langle\sstar\rangle\le C(c_I)\) on \(U\)
  & Proved mod (R) & \PROVEDMOD{saturation} \\
\(\sstar\) transfer: lower semicontinuity
  & \textbf{Proved} (the available direction) & \PROVED \\
\(\sstar\) transfer: equality
  & Proved mod (R),(N) & \PROVEDMOD{N} \\
Pull-back of H1 to the sequence
  & \textbf{Forbidden} (no reverse inequality) & Forbidden \\
\texttt{prop:localized-cf} on profiles
  & \textbf{Fails} (\(\Omega=\infty\)) & \IMPOSSIBLE{} as written \\
Local alignment from \(\sstar=0\)
  & \textbf{Proved} & \PROVED \\
Global alignment upgrade \textbf{(A-up)}
  & \textbf{Open} --- largest new gap & \OPEN{} (unchanged) \\
Aligned \(\Rightarrow\) 2D (\(\partial_3U=0\))
  & \textbf{Proved} mod (R) & \PROVED \\
2D bounded ancient Liouville
  & Cited \cite{KNSS2009}; class caveat & cited, caveat stands \\
Aligned horn
  & Proved mod (R),(A-up) & \PROVEDMOD{A-up} \\
Naive finite-budget argument
  & \textbf{Proved impossible} & \IMPOSSIBLE \\
Log-time amplitude ODE
  & Proved mod (R), attainment & \PROVEDMOD{attainment} \\
Averaged necessary condition
  & \textbf{Proved} mod (R), attainment & \PROVEDMOD{attainment} \\
Disorder--depletion target
  & \textbf{Open target}, marginal & \OPEN{} (unchanged) \\
\quad (T4a) average-to-point for \(w\)
  & Open; plugs into the \(K\)-family & \OPEN \\
\quad (T4b) profile depletion of \(\alpha\)
  & Open; needs a new far field & \OPEN \\
Two-step reduction
  & Proved mod (R),(N\('\)) & \PROVEDMOD{N\('\)} \\
Reframed dichotomy
  & Conditional on 5 inputs & \CONDITIONAL{4 inputs} \\
Type-II exclusion
  & Open (outside all present methods) & \OPEN{} (unchanged) \\
\end{longtable}}

\section{The ceiling, stated plainly}
\label{sec:ceiling}

\begin{record}
\textbf{\texttt{rem:s39-distance}} (line~7703), the document's own words.
Suppose both horns closed completely --- \textbf{(A-up)} proved and
\texttt{target:disorder-depletion} proved, with \textbf{(R)}, \textbf{(N)},
\textbf{(N\('\))} discharged. \emph{The conclusion would be: no Type-I
singularity. That is strictly less than the Clay problem.} What would remain:

\smallskip
\textbf{(1) Type-II exclusion.} Every estimate in \S33--\S39 uses the Type-I
rates, and the profile extraction uses them \emph{essentially}: without
\(M(t)\lesssim(T-t)^{-1}\) the rescaled family \eqref{eq:rescale} has no
uniform bound and there is no limit object at all. ``A Type-II singularity
would rescale to nothing, and the entire reframe would be vacuous against
it.'' This is outside these methods, and no part of the document should be
read as progress on it.

\smallskip
\textbf{(2) The nontriviality input.} Even in the Type-I case, \textbf{(N)} is
not a technicality. Until it is discharged, ``no Type-I singularity'' should
be read as ``no Type-I singularity \emph{at a non-degenerate point}''.

\smallskip
\textbf{(3) Marginality of the surviving target.}
\texttt{target:disorder-depletion} must beat the constant \(1\) exactly, and
the programme's history is that marginal targets are where errors hide.
\end{record}

Three things about that remark deserve emphasis, because a reader skimming
this book for the good news will otherwise mis-weight them.

First, it is not a caveat appended to a result. It is a statement about the
\emph{construction}, and the reason is structural rather than technical: the
compactness that makes the whole reframe possible is powered by the Type-I
rates, so the machinery cannot see the regime where those rates fail. A
sharper estimate does not help. A different tool is required, and none is on
offer here.

Second, it has survived every subsequent session unchanged. \Sn{41},
\Sn{44} and \Sn{48} all record explicitly that they do not touch it. The
\Sn{48} residue list says it ``stands verbatim''; the \Sn{48} scope remark
says the session ``does not bring the programme nearer to the Clay problem
than \texttt{rem:s39-distance} already states''. Wall~3
(Chapter~\ref{ch:walls}) is this remark, promoted to a wall.

Third, the honest one-line summary of \Sn{39} is the one the document itself
gives, and it is not a claim of progress toward the Prize: \emph{the single
deep gap of \Sn{38} has been replaced by two structurally unrelated gaps, one
of which is now complete except for a single clearly stated propagation
hypothesis, and one of which has yielded a proved averaged necessary condition
and a sharply-posed marginal target.} ``No claim is made beyond this.''

\begin{obsv}[What the reframe bought, and what it cost]
\label{obs:reframe-ledger}
\emph{Bought:} a change of object that made the aligned side nearly complete;
a genuinely new lemma (\texttt{lem:harmonic-reduction}) supplying the
dimensional reduction without which no Liouville theorem applies; a proof that
the obvious infinite-past argument is impossible, which saved later sessions
from attempting it; a log-time monotone quantity, which is the correct
framework the impossibility proof pointed to; and --- as a consequence
nine sessions later --- a branch on which the Calder\'on--Zygmund logarithm
turned out never to be incurred.

\smallskip
\emph{Cost:} one new hypothesis larger than anything it removed
(\textbf{(A-up)}); a proof that the \Sn{37} depletion machinery is unavailable
on profiles, invalidating the natural attack on the hard half of the remaining
target; and a ceiling, stated in the same section, that no amount of further
work inside the construction can raise.
\end{obsv}
